% FIGURE NEEDED: ray diagram -- two distant objects A and B whose rays cross
% at angle $\theta$ at a circular aperture of diameter $D$ (representing the
% eye); A' and B' are the central maxima of their diffracted images on the far
% side (2007 theory solutions, problem 1.13).

A and B are two distant objects. $A'$ and $B'$ are the central maxima of
their diffracted images. The angular position $\phi_1$ of the first minimum
relative to central maximum of each diffraction pattern is given by
$\phi_1 = 1.22\,\dfrac{\lambda}{D}$.

According to Lord Rayleigh the minimum angle of resolution is $\phi_1$.
Hence the images of A and B will be resolved if $\theta > \phi_1$,
$D > 1.22\,\dfrac{\lambda}{\theta}$.

We may take distance AB to be the diameter of the crater. Hence the
diameter $D$ of the eye's aperture must be, at least,
$D_{\min} = 1.22\,\dfrac{\lambda}{\theta}$.
\begin{align*}
\theta &= \frac{80\ \mathrm{km}}{\text{distance from Earth to Moon}}\\
&= \frac{80 \times 10^{3}\ \mathrm{m}}{3.844 \times 10^{8}\ \mathrm{m}}
 = 2.081 \times 10^{-4}\ \text{radian}
\end{align*}
\[ \lambda = 500 \times 10^{-9}\ \mathrm{m} \]
\begin{align*}
\therefore\ D_{\min}
&= \frac{1.22 \times 500 \times 10^{-9}}{2.081 \times 10^{-4}}\ \mathrm{m}
 = 2.93 \times 10^{-3}\ \mathrm{m}\\
&= 2.9\ \mathrm{mm}
\end{align*}
Hence, it is possible to resolve the 80\,km-diameter crater with naked eye.
